\documentclass[10pt,a4paper]{amsart}
\usepackage[margin=19mm]{geometry}
\usepackage{amsmath,amssymb,mathtools}
\usepackage[scr=rsfs,cal=euler]{mathalfa}
\usepackage{xfrac}
\usepackage[unicode,hidelinks,bookmarks=false]{hyperref}
\newcommand{\ie}{i.e.\ }
\newcommand{\cf}{cf.\ }
\newcommand{\resp}{respectively\ }
\newcommand{\Subsection}{\subsection}
\providecommand{\nobreakdash}{\nobreak\hbox{-}}
\setlength{\emergencystretch}{2em}
\multlinegap0pt
\usepackage{bm}
\usepackage{bbm}
\usepackage{dsfont}
\usepackage{xspace}
\usepackage{cancel}
\usepackage{mathtools}
\usepackage{amsthm}
\makeatletter
\@ifundefined{thm}{\newtheorem{thm}{Theorem}}{}
\@ifundefined{defn}{\newtheorem{defn}[thm]{Definition}}{}
\makeatother
\title[A new definition of upward planar order]{A new definition of upward planar order}
\author{Ting Li, Xuexing Lu}
\date{Source: arXiv:2308.06992v4 (CC BY 4.0)}
\begin{document}
\maketitle

\noindent\emph{Source and licence.} A new definition of upward planar order. Version arXiv:2308.06992v4 (CC BY 4.0), distributed under the Creative Commons Attribution 4.0 International licence, as recorded in the arXiv licence metadata for this exact version and shown on the abstract page https://arxiv.org/abs/2308.06992v4. Full rights notice: licenses/arxiv-2308.06992-CC-BY-4.0.txt.\par\smallskip

\noindent\textit{Editorial scope.} Counted unit: the Thm whose source label is \texttt{eq} in the article (source0.tex lines 319--321, source label \texttt{eq}), delivered together with the article's own complete original proof of it (source0.tex lines 322--336, one \texttt{proof} environment of 15 lines). No mathematics is written, repaired or re-derived: statement and proof are the article's own text line for line. Required context retained uncounted: upo (lines 64--72); upo1 (lines 231--244). Every definition, display and label of those lines is retained unchanged. Omitted from this package and not redistributed: 1--63, 73--230, 245--318, 337--373. Nothing inside a retained excerpt is omitted or altered: every retained body is delivered exactly as the article's lines are written, so no omission receipt of any kind accompanies this package. Citations: the retained excerpts carry no citation command at all, so no bibliography entry of the article is transcribed and no reference is redistributed. Reference adaptation: none was needed and none was made; every reference inside the delivered unit resolves inside it, so no unresolved reference or citation is produced. Printed slips kept literal: no printed word, symbol, hypothesis or number of the counted statement or of its proof was altered. Layout and class: the article's own class is \texttt{article} with packages including geometry, fontenc, inputenc, authblk, graphics, tikz, pgf, subfigure, here, caption, float, babel, xy, marvosym; the wrapper is the lane's pinned \texttt{amsart} editorial wrapper at 10pt on A4, which restates the theorem-like environments the retained text needs and adds the article's own macro definitions that the retained text needs (none), with extra wrapper packages (bm, bbm, dsfont, xspace, cancel, mathtools). The delivered numbering is the unit's own. Assets: no retained span contains a figure environment, a graphics-inclusion command, a PDF-page inclusion, a table or a TikZ picture; no image file is copied into this package, the package declares no asset dependency and no image file is redistributed with it. Licence: version arXiv:2308.06992v4 of this article is distributed under the Creative Commons Attribution 4.0 International licence, as recorded in the arXiv licence metadata for this exact version and shown on the abstract page https://arxiv.org/abs/2308.06992v4. A full rights notice is delivered at licenses/arxiv-2308.06992-CC-BY-4.0.txt. Characters: every retained excerpt is pure ASCII, so the delivered engine, XeLaTeX with its default Latin Modern text font, reports no missing character.
\begin{defn}\label{upo}
An \textbf{upward planar order} on acyclic directed graph $G$ is a linear order $\prec$ on the edge set $E(G)$, such that

$(U1)$ $e_1\to e_2$ implies that $e_1\prec e_2$;

$(U2)$ for any vertex $v$, $\overline{I(v)}\cap \overline{O(v)}=\emptyset$ and $\overline{E(v)}=\overline{I(v)}\sqcup \overline{O(v)}$;

$(U3)$ for any two  vertices $v_1$ and $v_2$, $I(v_1)\cap \overline{I(v_2)}\neq\emptyset$ implies that $\overline{I(v_1)}\subseteq \overline{I(v_2)}$, and $O(v_1)\cap \overline{O(v_2)}\neq\emptyset$ implies that $\overline{O(v_1)}\subseteq \overline{O(v_2)}$.
\end{defn}

\begin{defn}\label{upo1}
An \textbf{upward planar order} on acyclic directed graph $G$ is a linear order $\prec$ on the edge set $E(G)$, such that

$(Q1)$ $e_1\rightarrow e_2$ implies that $e_1\prec e_2$;

$(Q2)$ if $e_1\prec e\prec e_2$ and $e_1$, $e_2$ are adjacent, then
\begin{equation*}
\begin{cases}
I(t(e))\subseteq \overline{I(v)},&  \text{if}\ t(e_1)=t(e_2)=v;\\
O(s(e))\subseteq \overline{O(v)},&  \text{if}\ s(e_1)=s(e_2)=v;\\
\text{either}\ I(t(e))\subseteq \overline{I(v)}\  \text{or}\ O(s(e))\subseteq \overline{O(v)},&  \text{if}\ t(e_1)=s(e_2)=v.
\end{cases}
\end{equation*}
\end{defn}

\begin{thm}\label{eq}
Definition \ref{upo} and Definition \ref{upo1} are equivalent.
\end{thm}

\begin{proof}
Clearly, $(U1)=(Q1)$. We only need to show that $(U2)+(U3)\Longleftrightarrow (Q2)$.

$(\Longleftarrow)$. To show $(Q2)\Longrightarrow(U2)$, we only need to show that for any processive vertex $v$, $O(v)^-=I(v)^++1$. We prove this by contradiction. Suppose there exist a processive vertex $v$ and an edge $e$ such that $I(v)^+\prec e \prec O(v)^-$. Clearly, $v=t(I(v)^+)=s(O(v)^-)$, then by $(Q2)$, we have either $e\in I(t(e))\subseteq \overline{I(v)}$ or $e\in O(s(e))\subseteq\overline{O(v)}$, both of which  contradict $I(v)^+\prec e \prec O(v)^-$.

Now we show $(Q2)\Longrightarrow(U3)$. Assume $I(v_1)\cap \overline{I(v_2)}\neq \emptyset$, then there exists $e\in I(v_1)$ such that $I(v_2)^-\prec e \prec I(v_2)^+$. Clearly, $v_2=t(I(v_2)^-)=t(I(v_2)^+)$, then by $(Q2)$, we have $I(v_1)=I(t(e))\subseteq \overline{I(v_2)}$. Similarly, $O(v_1)\cap \overline{O(v_2)}\neq \emptyset$ means that there exists $e\in O(v_1)$ such that $O(v_2)^-\prec e \prec O(v_2)^+$. By $v_2=s(O(v_2)^-)=s(O(v_2)^+)$ and $(Q2)$, we have $O(v_1)=O(s(e))\subseteq \overline{O(v_2)}$.

$(\Longrightarrow)$. Assume $e_1\prec e\prec e_2$ and $e_1$, $e_2$ are adjacent, then there are three cases.

\noindent\textbf{Case 1:} $t(e_1)=t(e_2)=v$. In this case, $e\in I(t(e))\cap\overline{I(v)}$, by $(U3)$, we have $I(t(e))\subseteq\overline{I(v)}$.

\noindent\textbf{Case 2:} $s(e_1)=s(e_2)=v$. In this case, $e\in O(s(e))\cap\overline{O(v)}$, by $(U3)$, we have $I(s(e))\subseteq\overline{O(v)}$.

\noindent\textbf{Case 3:} $t(e_1)=s(e_2)=v$. In this case, $e\in \overline{E(v)}$, then by $(U2)$, we have either $e\in \overline{I(v)}$ or $e\in \overline{O(v)}$, which reduces the proof to case 1 or case 2, respectively.
\end{proof}

\end{document}
